Let $G=\langle a_1,\dots,a_r\rangle$ be a free group. A word on the
  letters $L=\{a_1,a_1^{-1},\dots,a_r,a_r^{-1}\}$ is called reducible
  if it equals $1\in G$. A word is called quasi-reduced if it does not
  contain three consecutive letters $\ell\ell^{-1}\ell$ for any
  $\ell \in L$.

  Fix a reducible quasi-reduced word $w=w_1\cdots w_{2n}$ (with $w_i\in L$). A
  matching $M$ consists of $n$ intervals, each properly embedded in
  $\RR\times\RR_+$, satisfying the following:
  \begin{itemize}[leftmargin=*]
  \item Any two intervals intersect
    transversally, and there are no triple points.
  \item $\bdy M=\{1,\dots,2n\}\times 0\subset\RR\times 0\subset
    \RR\times\RR_+$.
  \item For each component $R$ of
    $\RR\times\RR_+\setminus M$, and for each component $A$ of $M\cap \bdy\overline{R}$, $\bdy A$ consists of some two points $(i,0)$ and $(j,0)$ with $w_i=w_j^{-1}$. 
  \end{itemize}
  A $k$-crossing
  matching is a matching where the total number of self-intersections
  is $k$.

  Let $F_w$ be the following CW complex whose $k$-cells correspond to
  isotopy classes of $k$-crossing matchings. Given a $k$-crossing
  matching $M$ with crossing set $X(M)$, the associated $k$-cell is
  given by $C(M)\defeq \prod_{c\in X(M)}I_c$, where $I_c$ is an
  interval whose endpoints correspond to the two ways of locally
  resolving $M$ near $c$. The attaching maps are given as follows: On
  a facet of the form
  \[
    a\times\prod_{d\in X(M)\setminus\{c\}}I_{d},\qquad c\in X(M),a\in \bdy I_c,
  \]
  let $M_a$ be the $(k-1)$-crossing matching obtained by resolving $M$
  at $c$ according to $a$. Then the $(k-1)$-cell associated to $M_a$
  is given by
  $C(M_a)=\prod_{d\in X(M_a)}I_d=\prod_{d\in X(M)\setminus\{c\}}I_{d}$, and the
  above facet maps to this $(k-1)$-cell canonically.

  Is $F_w$ contractible?